\subsection{N-020: the abstract sign fibre, its radical, and an actual cyclic cell}
\label{sec:nonstandard-sign-fibre}

\paragraph{The three objects and the two traces.}
Put $S=\mathbb Q[D]$ and let $B_S$ be the actual polynomial order
of N-015, with the complete actual basis and table of N-019.
The sign fibre studied here is
\[
 \mathcal A=B_S/DB_S,\qquad \dim_{\mathbb Q}\mathcal A=114.
 \tag{N-104}
\]
This is an abstract algebra fibre.  It is distinct from its image
$\mathcal I$ in the tensor square of the Hecke algebra specialized
at $v=i$, and from VIII's unspecified $E(4)$-labelled table.
The latter specialization has $q=-1$ and $D=0$; extending (N-104)
to $\mathbb Q(i)$ gives the corresponding rational-coefficient
Laurent-order fibre.  The image and kernel below have rational
forms in the normalized opposite-KL coordinates of N-015.

Let $\tau_{\rm amb}$ be the pullback of the ambient Hecke
symmetrizing functional, defined on the standard basis by
$\tau(T_x\otimes T_y)=\mathbf1_{x=y=1}$.
Let $\tau_{\rm reg}(a)=\operatorname{Tr}_{\mathbb Q}(L_a)$ be
the trace of left multiplication on the abstract algebra
$\mathcal A$.  These are different functionals.  Exact rational
computation gives
\[
\begin{array}{c|r}
 \dim\mathcal A&114\\
 \dim\mathcal I&55\\
 \dim\ker(\mathcal A\longrightarrow\mathcal I)&59\\
 \operatorname{rank}(\tau_{\rm amb}(H_iH_j))&42\\
 \operatorname{rank}(\tau_{\rm reg}(H_iH_j))&5.
\end{array}
\tag{N-105}
\]
Thus the nullspace of the ambient pairing has dimension72;
it is not the Jacobson radical computed below.  The ambient image
recovers N-008, while (N-104) and the regular trace are new
abstract-fibre calculations from the actual free order.

\paragraph{An explicit split quotient.}
There is a surjective algebra homomorphism
\[
 \begin{gathered}
 \Phi=(\chi_0,\rho):\mathcal A\longrightarrow
                \mathbb Q\oplus M_2(\mathbb Q),\\
 \Phi(P_1)=\Phi(P_3)=(0,E_{12}),\qquad
 \Phi(P_2)=(0,4E_{21}).
 \end{gathered}
 \tag{N-106}
\]
Here $\chi_0$ is the character sending every $P_i$ to zero.
To see that the two-dimensional representation extends over the
actual order before specialization, use the polynomial model
\[
 P_1=P_3=\begin{pmatrix}D&1\\0&0\end{pmatrix},
 \qquad
 P_2=\begin{pmatrix}0&0\\(D-2)^2&D\end{pmatrix}.
 \tag{N-107}
\]
It is conjugate over $\mathbb Q(D)$ to the certified $S'_{22}$
GT model of N-010 by $Q=\left(\begin{smallmatrix}1&1\\0&-D\end{smallmatrix}\right)$:
the GT matrices are $Q^{-1}P_iQ$.  Since every element of $B_S$
is a polynomial combination of words in these generators, the
model sends the actual order into $M_2(S)$.  Relations hold there
because they hold in its certified generic representation and $M_2(S)$ is
torsion free.  This proves (N-106) on the abstract fibre, without
dividing by $D$ after specialization.  The images of $P_1,P_2$
and their products span the matrix summand; subtracting its
identity from $\Phi(1)$ also gives the scalar summand.

The complete table proves the following identity of linear
functionals, checked on all114 actual basis vectors:
\[
       \tau_{\rm reg}=78\chi_0+18\operatorname{tr}\rho.
 \tag{N-108}
\]
The right side takes value $78+18\cdot2=114$ on the identity.
All12996 product identities for $\Phi$ are also checked exactly.
In particular the regular pairing is the pullback of
$78ab+18\operatorname{tr}(XY)$ on
$\mathbb Q\oplus M_2(\mathbb Q)$, a nondegenerate form.
Its kernel is exactly $\ker\Phi$, of dimension109.

\paragraph{The kernel is precisely the Jacobson radical.}
This conclusion does not require confusing a symmetrizing
functional with a regular trace.  If $x\in\ker\Phi$, then
$x^m\in\ker\Phi$ for every $m\ge1$, and (N-108) gives
$\operatorname{Tr}(L_x^m)=\tau_{\rm reg}(x^m)=0$.
Newton's identities over $\mathbb Q$, applied up to degree114,
imply that the characteristic polynomial of $L_x$ is $t^{114}$.
Cayley--Hamilton gives $L_x^{114}=0$, hence $x^{114}=0$ by
applying to the identity.  Thus every element of the ideal
$\ker\Phi$ is nilpotent.  For any $a\in\mathcal A$, $ax$
lies in the same ideal, so $1-ax$ is invertible by a finite
geometric series.  The usual defining criterion for the
Jacobson radical gives $\ker\Phi\subseteq J(\mathcal A)$.
Conversely, its image in the split semisimple quotient in
(N-106) is zero, giving the reverse inclusion.  Therefore
\[
 J(\mathcal A)=\ker\Phi,\qquad
 \dim J(\mathcal A)=109,\qquad
 \mathcal A/J(\mathcal A)\simeq\mathbb Q\oplus M_2(\mathbb Q).
 \tag{N-109}
\]

The exact ambient and quotient matrices satisfy
$\operatorname{rank}[\mathrm{ambient}\mid\Phi]=55$.
Together with ambient rank55 this means that every quotient
coordinate factors linearly through the ambient image.
Hence $\ker(\mathcal A\to\mathcal I)\subseteq\ker\Phi$.
It follows that the ambient image has radical dimension50
and the same split semisimple quotient of dimension5.
This is an assertion about $\mathcal I$ as a quotient algebra;
it is not an identification of its dimension with (N-104).

\paragraph{The parabolic rank-three algebra survives abstractly.}
Write $a=P_1$ and $b=P_2$.  The ten words
\[
 1,\ a,\ b,\ ab,\ ba,\ aba,\ bab,\ abab,\ baba,\ ababa
 \tag{N-110}
\]
are independent in $\mathcal A$, and their span is closed under
$a,b$.  Thus they give the entire ten-dimensional parabolic
subalgebra.  This is also the abstract rank-three fibre, not
merely its seven-dimensional ambient image.

For completeness, the actual rank-three polynomial order is
free on (N-110) before specialization.  This presentation and
freeness are already known: GCT IV, Theorem11.16/equation(112),
printed p.~65; Blasiak (2015), Theorem3.7 and Corollary3.8,
preprint pp.~13--14.  The following reductions also check them
directly in the present convention.  Put
$L=(D-2)^2$, $u=D+L$, and $w=DL$.
The faithful rank-three generic models in N-010 give
$a^2=Da$, $b^2=Db$ and
\[
 ababa-babab=u(aba-bab)-w(a-b).
\]
Multiplication on the right by $b$, followed by these relations,
gives
\[
 ababab=Dababa+uabab-Duaba-wab+Dwa.
\]
The swapped relation handles $bababa$.  These monic reductions
span every word by (N-110).  The computed independence after
specialization in the free rank-four order implies their generic
independence, and hence independence over $S$.
Its abstract fibre therefore has dimension10.  The computed
rank10 inclusion in $\mathcal A$ proves that it stays injective
there.  Its ambient image has dimension7 and ambient pairing
rank6, as in the supplied sign-point computation and N-008.

\paragraph{A corrected cell is an actual matrix subalgebra.}
At $D=0$ set
\[
 \begin{gathered}
 \gamma_1=aba,\quad\gamma_{12}=abab,\quad
 \gamma_{21}=baba,\quad\gamma_2=bab,\\
 \Sigma=ababa-4aba=babab-4bab,\\
 h_1=\gamma_1+\tfrac14\Sigma=\tfrac14ababa,\qquad
 h_2=\gamma_2+\tfrac14\Sigma=\tfrac14babab,\\
 h_{12}=\gamma_{12},\qquad h_{21}=\gamma_{21}.
 \end{gathered}
 \tag{N-111}
\]
The element $\Sigma$ is nonzero in the abstract fibre,
$\Sigma^2=0$, and it annihilates the four $\gamma$'s on both
sides.  It vanishes in the ambient image.  These statements
are exact word identities in the table.  The uncorrected
four-vector span is not a subalgebra in the abstract fibre:
four of its products have the nonzero $4\Sigma$ term described
in the supplied E4 note, section6b.  Thus that term cannot be
dropped when answering the abstract-fibre question.

The correction in (N-111) removes it.  Under $\rho$, in the
displayed order, these four vectors are
$4E_{12},16E_{11},16E_{22},16E_{21}$, and their full table is
\[
\begin{array}{c|rrrr}
 &h_1&h_{12}&h_{21}&h_2\\\hline
 h_1&0&0&16h_1&4h_{12}\\
 h_{12}&16h_1&16h_{12}&0&0\\
 h_{21}&0&0&16h_{21}&16h_2\\
 h_2&4h_{21}&16h_2&0&0
\end{array}
\tag{N-112}
\]
Consequently their span $\mathcal C$ is an actual copy of
$M_2(\mathbb Q)$, with local identity
$e=(h_{12}+h_{21})/16$.  The element $f=1-e$ is an idempotent
annihilating $\mathcal C$ on both sides.  Hence
$\mathbb Qf\oplus\mathcal C$ is an explicit split semisimple
subalgebra of $\mathcal A$, and $\Phi$ restricts to an
isomorphism from it to the quotient (N-106).
This constructs a semisimple section using actual parabolic
words.  No assertion that $\mathcal C$ is an ideal of all of
$\mathcal A$ is needed.

\paragraph{The cyclic symmetry and its exact limit.}
The ambient trace pairing on the $h$'s is
\[
 G=\begin{pmatrix}
 0&0&0&64\\0&256&0&0\\0&0&256&0\\64&0&0&0
 \end{pmatrix}.
 \tag{N-113}
\]
Thus the pairing permutation is
$\pi=(h_1\ h_2)$, fixing $h_{12},h_{21}$, with weights
$(64,256,256,64)$.  If $c(i,j;k)$ denotes a coefficient in
(N-112), cyclicity of the trace, or direct verification of
all64 triples, gives
\[
 c(i,j;k)w_k=c(j,\pi k;\pi i)w_{\pi i}.
 \tag{N-114}
\]
The slot permutation $R(i,j,k)=(j,\pi k,\pi i)$ has order3.
The parabolic generator swap $a\leftrightarrow b$ induces
$\sigma=(h_1\ h_2)(h_{12}\ h_{21})$ on this cell.
Acting in all three slots, it commutes with $R$ and gives
$C_3\times C_2\simeq C_6$.  Word reversal induces the
anti-automorphism $t=(h_{12}\ h_{21})$; its action on slots
is $(i,j,k)\mapsto(tj,ti,tk)$ and conjugates $R$ to $R^{-1}$.
Together with $R$, this reflection gives $D_3$ of order6.
These are exact weighted coefficient-tensor symmetries.
The parabolic swap is not being identified with the full
rank-four Dynkin reversal $P_1\leftrightarrow P_3$.

The same permutation and weighted cyclic identity cannot
extend to a full basis of $\mathcal A$ containing $1$ and
these four cell labels, with all weights nonzero.
Indeed take $i=1$, $j=k=h_{12}$ in (N-114).
The left side is $w_{h_{12}}\ne0$.  The right side asks for
the coefficient of $\pi(1)$ in $h_{12}^2=16h_{12}$.
Since $\pi$ fixes $h_{12}$ and is a permutation,
$\pi(1)\ne h_{12}$; this coefficient is zero.
This is the same unit obstruction as N-001, now inside the
actual abstract rank-four fibre.  Separately, rank42 in
(N-105) excludes any full114-element nonzero-weight monomial
Gram matrix for this ambient functional, irrespective of
basis.  Neither statement denies ordinary trace cyclicity
or every possible construction using a different functional
and different labels.

\paragraph{Certificate and source scope.}
The constructors are \nolinkurl{sign_fibre_abstract.py} and
\nolinkurl{sign_fibre_quotient.py}.  They specialize every
coefficient of N-019's complete polynomial table, compute
the114 multiplication traces, and store both full Gram
matrices.  The ambient matrix is obtained from the actual
300 normalized coordinates: coordinates with an odd power
of $\delta=v+v^{-1}$ vanish at $v=i$, and the others are
evaluated at $D=0$.  The tensor symmetrizing functional is
evaluated in the certified opposite-KL basis with the
off-diagonal factor two retained.

The separate verifier \nolinkurl{verify_sign_fibre.py}
uses no constructor imports.  It checks the stored fibre
table against N-019, the generator preimages against the
generic GT matrices, every quotient product, the trace
identity (N-108), exact ranks, all parabolic closure and
cell identities, and the full64-triple symmetries.
Its finite identities prove statements about the whole
algebra because the actual114-element polynomial basis
and all products have already been certified; no word
degree cutoff is used to infer full closure.

The inherited rank-three input is the supplied
\nolinkurl{OURS_CONSULT_a157422a_COUNTING_CONSTRUCTION.md},
Theorem2, together with
\nolinkurl{OURS_E4_RANK_FOUR_20261121.md}, section6b and its
\nolinkurl{gram_signpoint.py}/\nolinkurl{cell_c3_check.py}
locators.  Their trace ranks and weighted cell permutation
are retained, with the abstract $\Sigma$ term now handled
explicitly.  The radical criterion and matrix-algebra
arguments are standard finite-dimensional algebra; the
particular114-dimensional fibre, its exact ranks, and
the explicit corrected section are the computations here.
GCT IV, Remark11.9, printed p.~61, already distinguishes
abstract and ambient fibres and identifies the radical at
its ordinary parameter specialization $v=1$, or $D=4$.
It is not the present sign fibre.  The bounded primary
literature audit is \nolinkurl{N020_literature_scope.md}.
This answers the indicated local-cell and full-extension
question for the actual rational order.  It does not supply
a missing $E(4)$ basis or a canonical positive dimension ring.
